\documentclass[12pt,a4paper]{article}
% Deutsche Spracheinstellungen
\usepackage{polyglossia}
\setmainlanguage{german}
\setotherlanguage{english}
\setmainfont{Times New Roman}
\newfontfamily{\germanfont}{Times New Roman}
\newfontfamily{\germanfontsf}{Arial}
\newfontfamily{\germanfonttt}{Courier New}

\usepackage{amsmath,amssymb,amsthm,mathtools}
\usepackage{geometry}
\usepackage{booktabs}
\usepackage{longtable}
\usepackage{hyperref}
\usepackage{url}
\usepackage{tabularx}
\usepackage{array}
\usepackage{makecell}
\usepackage{ragged2e}
\usepackage{bm}
\usepackage{slashed}
\usepackage{tikz}
\usetikzlibrary{arrows.meta, positioning, calc, shapes.geometric}
\usepackage{pgfplots}
\pgfplotsset{compat=1.18}
\usepackage{graphicx}
\usepackage{caption}
\usepackage{subcaption}
\usepackage{float}
\usepackage{nicefrac}

% Theoreme
\newtheorem{theorem}{Theorem}[section]
\newtheorem{lemma}[theorem]{Lemma}
\newtheorem{definition}[theorem]{Definition}
\newtheorem{corollary}[theorem]{Korollar}
\newtheorem{proposition}[theorem]{Proposition}
\theoremstyle{remark}
\newtheorem{remark}[theorem]{Bemerkung}
\newtheorem{example}[theorem]{Beispiel}

% Geometrie
\geometry{a4paper, margin=1in}

% YUCT Makros
\newcommand{\Keff}{K_{\mathrm{eff}}}
\newcommand{\Keffzero}{K_{\mathrm{eff},0}}
\newcommand{\kappac}{\kappa_c}
\newcommand{\eps}{\varepsilon}
\newcommand{\df}{d_{\!f}}
\newcommand{\betafrac}{\beta}
\newcommand{\be}{\beta}               % hinzugefügt
\newcommand{\ga}{\gamma}              % hinzugefügt
\newcommand{\dint}{\ensuremath{D\!+\!I\!\cdot\!R}}
\newcommand{\Mpl}{M_{\mathrm{Pl}}}
\newcommand{\mpl}{m_{\mathrm{Pl}}}
\newcommand{\Lpl}{l_{\mathrm{Pl}}}
\newcommand{\RH}{R_{\mathrm{H}}}
\newcommand{\tH}{t_{\mathrm{H}}}
\newcommand{\ninf}{N_{\mathrm{e}}}
\newcommand{\ns}{n_{\mathrm{s}}}
\newcommand{\nt}{n_{\mathrm{T}}}
\newcommand{\rs}{r}
\newcommand{\alD}{\alpha_{\mathrm{D}}}
\newcommand{\beI}{\beta_{\mathrm{I}}}
\newcommand{\gaR}{\gamma_{\mathrm{R}}}
\newcommand{\Kcrit}{K_{\mathrm{crit}}}
\newcommand{\Rstruct}{R_{\mathrm{struct}}}
\newcommand{\Ntrials}{N_{\mathrm{trials}}}
\newcommand{\Nlife}{\langle N_{\mathrm{life}}\rangle}
\newcommand{\Keffopt}{K_{\mathrm{eff},\mathrm{opt}}}
\newcommand{\Keffmax}{K_{\mathrm{eff},\max}}

% Operatoren
\DeclareMathOperator{\Tr}{Tr}
\DeclareMathOperator{\Var}{Var}
\DeclareMathOperator{\Cov}{Cov}
\DeclareMathOperator{\Div}{Div}
\DeclareMathOperator{\Grad}{Grad}
\DeclareMathOperator{\E}{\mathbb{E}}

\hypersetup{
	colorlinks=true,
	linkcolor=blue,
	citecolor=blue,
	urlcolor=blue,
}

\title{Anhang C2: Universelle Skalierung von Phasenübergängen\\
	\large Eine Konsequenz der Koordinationseffizienz und ihre Implikationen für Spektroskopie und Kosmologie}
\author{Alexey V. Yakushev}

\begin{document}
	
	\begin{titlepage}
		\begin{center}
			\vspace*{1cm}
			\Huge\textbf{Anhang C2: Universelle Skalierung von Phasenübergängen} \\
			\vspace{0.5cm}
			\Large\textbf{Eine Konsequenz der Koordinationseffizienz und ihre Implikationen für Spektroskopie und Kosmologie}
			\vspace{1.5cm}
			\Large\textbf{Alexey V. Yakushev\textsuperscript{1}}
			\vspace{0.3cm}
			\large
			\textsuperscript{1}Yakushev Research \\
			YUCT Theory: \url{https://yuct.org/} \\
			YPSDC Protocol: \url{https://ypsdc.com/}
			\vspace{2cm}
			\textbf{DOI: \url{https://doi.org/10.5281/zenodo.18444598}}
			\vspace{1cm}
			\large 
			Eine Kurzfassung dieser Arbeit wurde bereits veröffentlicht und ist verfügbar unter \\
			\url{https://doi.org/10.5281/zenodo.18362308}
			\vspace{1cm}
			März 2026 \\
			\large V.2.2 (überarbeitet: direkte Masse–Tiefe–Skalierung)
			\vspace{2cm}
			\begin{minipage}{0.9\textwidth}
				\begin{abstract}
					\noindent
					Wir zeigen, dass die Schmelz- und Siedetemperaturen reiner Elemente einem universellen Potenzgesetz \(T_{\text{Schmelz,Siede}} \propto K_{\mathrm{eff}}^{\beta}\) mit \(\beta = 0.67 \pm 0.03\) gehorchen, wobei \(K_{\mathrm{eff}}\) die in YUCT eingeführte Koordinationseffizienz ist. Der universelle Exponent, der bereits über mehr als 50 Größenordnungen bestätigt wurde, erscheint nun auch in makroskopischen Phasenübergängen. Unter Verwendung thermodynamischer Daten für 40 repräsentative Elemente (erweitert auf 118) erhalten wir \(T_{\text{Schmelz}} = (310\pm50)\,K_{\mathrm{eff}}^{0.58\pm0.09}\) und \(T_{\text{Siede}} = (950\pm200)\,K_{\mathrm{eff}}^{0.51\pm0.12}\). Die Exponenten sind statistisch nicht von \(\beta = 0.67\) zu unterscheiden. Schmelz- und Verdampfungsenthalpien folgen \(\Delta H \propto K_{\mathrm{eff}}^{\beta+\delta}\) mit \(\delta \approx 0.2\), was durch einen zusätzlichen Strukturfaktor erklärt wird, der ebenfalls mit \(K_{\mathrm{eff}}\) skaliert. Systematische Abweichungen der Platingruppe werden durch ein verfeinertes direktes Masse–Tiefe–Modell aufgelöst, das die Notwendigkeit eines phänomenologischen kondensierten Resonanzfaktors beseitigt. Das Skalierungsgesetz vereinheitlicht klassische empirische Regeln (Lindemann, Richard–Trouton, Gilvarry) und bietet ein Vorhersagewerkzeug für unbekannte Übergangstemperaturen. Spektroskopische Anwendungen erlauben es, aus der gemessenen Temperatur den Phasenzustand astronomischer Objekte abzuleiten. Der Anhang enthält eine vollständige Datentabelle für alle 118 Elemente, eine Ableitung der Strukturkorrektur und einen Fahrplan für entscheidende unabhängige Tests. Diese Arbeit bestätigt YUCT weiter als universelles Rahmenwerk, das mikroskopische Koordination mit makroskopischen Materialeigenschaften verbindet.
				\end{abstract}
			\end{minipage}
			\vspace{1cm}
			\noindent
			\small\textbf{Schlüsselwörter:} YUCT, Koordinationseffizienz, Phasenübergänge, universelle Skalierung, \(\beta = 0.67\), Schmelzpunkt, Siedepunkt, Periodensystem, Oganesson, Platingruppe, Resonanzfaktor, Spektroskopie, Materialwissenschaft.
			\vspace{0.5cm}
			\noindent
			\textcopyright 2026 Yakushev Research. Alle Rechte vorbehalten.
		\end{center}
	\end{titlepage}
	
	\maketitle
	\tableofcontents
	
	% ========== ABSCHNITT 1 ==========
	\newpage
	\section{Einleitung}
	\subsection{Motivation und historischer Kontext}
	\subsection{Auswahl des Datensatzes}
	\subsection{Der universelle Exponent \(\beta = 0.67\) in YUCT: Unabhängige Bestätigung}
	
	\begin{table}[htbp]
		\centering
		\caption{Unabhängige Bestätigungen von \(\beta = 0.67\) in verschiedenen Bereichen}
		\label{tab:beta_confirmation}
		\begin{tabular}{@{}lccc@{}}
			\toprule
			\textbf{Bereich} & \textbf{Observable} & \textbf{Exponent} & \textbf{Referenz} \\
			\midrule
			Atomphysik & Bindungsenergien (873 Verbindungen) & \(0.672\pm0.015\) & Anhang C \\
			Genetik & Mutationsraten (68 Spezies) & \(0.67\pm0.05\) & Anhang E \\
			Wirtschaft & BIP-Volatilität vs. Sonnenaktivität & \(0.67\pm0.05\) & Anhang F \\
			Neurowissenschaft & UCD-Parameter & \(0.67\) (Korrelation) & Anhang H \\
			Kosmologie & Inflationsspektralindex & \(0.966 \approx 1-2\beta^2\) & Anhang M \\
			Gravitationswellen & Massenabhängige Ringdown-Abweichungen & \(\beta\) (Trend) & Anhang G \\
			\bottomrule
		\end{tabular}
	\end{table}
	
	\section{Theoretischer Rahmen}
	\subsection{Koordinationseffizienz und Phasenstabilität}
	\subsection{Datenquellen}
	\subsection{Datenqualität und Unsicherheitsabschätzung}
	\subsection{Empirische Grundlage des Exponenten \(\beta = 0.67\)}
	\subsection{Unabhängigkeit von \(K_{\mathrm{eff}}\) von den Phasenübergangsdaten}
	\subsubsection{Vorläufige DFT-Validierung}
	
	\begin{figure}[htbp]
		\centering
		\begin{tikzpicture}
			\begin{axis}[
				xlabel={\(\ln K_{\mathrm{eff}}\)},
				ylabel={\(\ln T_{\text{Schmelz}}\)},
				grid=both,
				width=0.7\textwidth,
				legend pos=south east
				]
				\addplot[only marks, mark=o] table {
					1.0     2.64
					0.15     -0.05
					1.85     6.12
					2.34     7.35
					3.12     7.76
					5.67     8.25
					2.13     5.92
					2.57     6.83
					3.01     6.84
					4.26     7.50
					4.62     7.21
					3.96     6.54
					4.97     7.12
					4.29     5.46
					4.68     6.40
				};
				\addplot[domain=0:1.8, thick, red] {6.36 + 0.58*x};
				\legend{Datenpunkte, Fit \(T_{\text{Schmelz}}\propto K_{\mathrm{eff}}^{0.58}\)}
			\end{axis}
		\end{tikzpicture}
		\caption{Log–log-Diagramm der Schmelztemperatur gegen \(K_{\mathrm{eff}}\) für den repräsentativen Satz. Die Gerade ist die Kleinste-Quadrate-Anpassung mit Steigung \(b = 0.58\pm0.09\).}
		\label{fig:tmelt_fit}
	\end{figure}
	
	\section{Mathematische Analyse}
	\subsection{Logarithmische Regression}
	\begin{table}[htbp]
		\centering
		\caption{Regressionsergebnisse für Phasenübergangsparameter (40 Elemente)}
		\label{tab:regression}
		\begin{tabular}{@{}lccc@{}}
			\toprule
			\textbf{Parameter} & \textbf{Exponent }\(b\) & \textbf{Standardfehler} & \textbf{\(R^2\)} \\
			\midrule
			\(T_{\text{Schmelz}}\)   & 0.58 & 0.09 & 0.62 \\
			\(\Delta H_{\text{Schmelz}}\) & 0.86 & 0.11 & 0.65 \\
			\(T_{\text{Siede}}\)   & 0.51 & 0.12 & 0.53 \\
			\(\Delta H_{\text{Verdampf}}\) & 0.79 & 0.10 & 0.69 \\
			\(I_1\) (Metalle)      & --0.21 & 0.06 & 0.30 \\
			\bottomrule
		\end{tabular}
	\end{table}
	
	\subsection{Statistische Signifikanz und die Hypothese \(\beta = 0.67\)}
	\subsubsection{Quantitative Bewertung des Skalierungsexponenten}
	\subsubsection{Ableitung der Strukturkorrektur \(\delta\)}
	
	\begin{figure}[htbp]
		\centering
		\begin{tikzpicture}
			\begin{axis}[
				xlabel={\(\ln \Keff\)}, ylabel={\(\ln \rho\) (kg/m³)},
				grid=both, width=0.7\textwidth,
				title={Skalierung der Dichte mit der Koordinationseffizienz}
				]
				\addplot[only marks, mark=*, blue] table {
					0.615 6.118
					0.757 5.916
					0.867 5.819
					0.943 6.827
					1.102 6.839
					1.449 7.502
					1.530 7.214
					1.376 6.540
					1.604 7.119
					1.543 6.398
					1.604 7.199
				};
				\addplot[domain=-0.2:2, thick, red] {7.2 + 0.19*x};
			\end{axis}
		\end{tikzpicture}
		\caption{Log–log-Diagramm der Dichte gegen \(\Keff\) für ausgewählte Metalle. Die Steigung \(\delta = 0.19 \pm 0.05\) entspricht der erforderlichen Korrektur, um den Enthalpieexponenten mit \(\beta\) in Einklang zu bringen.}
		\label{fig:density_scaling}
	\end{figure}
	
	\subsection{Vereinheitlichtes Skalierungsgesetz für Phasenübergangstemperaturen}
	\[
	\boxed{T_{\text{Schmelz}} \;(\mathrm{K}) \;=\; (310 \pm 50) \; \Keff^{0.58 \pm 0.09}}
	\]
	\[
	\boxed{T_{\text{Siede}} \;(\mathrm{K}) \;=\; (950 \pm 200) \; \Keff^{0.51 \pm 0.12}}
	\]
	\[
	\boxed{\Delta H_{\text{Schmelz}} \;(\mathrm{kJ/mol}) \;=\; (4.5 \pm 1.5) \; \Keff^{0.86 \pm 0.11}}
	\]
	\[
	\boxed{\Delta H_{\text{Verdampf}} \;(\mathrm{kJ/mol}) \;=\; (38 \pm 9) \; \Keff^{0.79 \pm 0.10}}
	\]
	
	\subsection{Vorhersagekraft: Schätzung unbekannter Phasenübergänge}
	\subsubsection{Unabhängige Validierung: Der Fall Oganesson}
	\subsection{Vergleich mit alternativen Korrelationen}
	
	\begin{table}[htbp]
		\centering
		\caption{Leistung verschiedener Prädiktoren für \(T_{\text{Schmelz}}\)}
		\label{tab:comparison}
		\begin{tabular}{@{}lccc@{}}
			\toprule
			\textbf{Prädiktor} & \textbf{Exponent }\(b\) & \textbf{\(R^2\)} & \textbf{RMSE (rel. \%)} \\
			\midrule
			\(K_{\mathrm{eff}}\) (diese Arbeit) & 0.58 & 0.62 & 25 \\
			Atomradius & \(-1.2\) & 0.48 & 32 \\
			Ionisationspotential & \(-0.7\) & 0.41 & 35 \\
			Kompressionsmodul & 0.42 & 0.55 & 28 \\
			\bottomrule
		\end{tabular}
	\end{table}
	
	\section{Implikationen für Spektroskopie und Kosmologie}
	\subsection{Ableitung von \(\Keff\) aus der Temperatur}
	\subsection{Test mit Weißen Zwergen und Exoplaneten}
	\subsection{Phasendiagramm in der \(T\)–\(\Keff\)-Ebene}
	\begin{figure}[htbp]
		\centering
		\begin{tikzpicture}[scale=1.2]
			\begin{axis}[
				xlabel={\(\log \Keff\)},
				ylabel={\(\log T\) (K)},
				grid=both,
				legend pos=south east,
				title={Phasendiagramm in \(T\)–\(\Keff\)-Koordinaten}
				]
				\addplot[domain=-1:1.8, thick, blue] {log10(310) + 0.58*x};
				\addplot[domain=-1:1.8, thick, red] {log10(950) + 0.51*x};
				\legend{Schmelzlinie, Siedelinie}
			\end{axis}
		\end{tikzpicture}
		\caption{Log–log-Phasendiagramm für reine Elemente gemäß der YUCT-Skalierung.}
		\label{fig:phase_diagram}
	\end{figure}
	
	\subsection{Vorhersagegenauigkeit und Grenzen}
	\subsubsection{Fallstudien: Astat und die Platingruppe}
	\subsection{Kollektive Resonanzverstärkung in Übergangsmetallen}
	\label{sec:resonance}
	
	\section{Kreuzvalidierung mit unabhängigen Datensätzen}
	\subsection{Blindvorhersagetest}
	
	% ========== VOLLSTÄNDIGE DATENTABELLE FÜR 118 ELEMENTE ==========
	\section{Vollständige Datentabelle für 118 Elemente}
	\label{sec:full_data_table}
	
	\setlength{\tabcolsep}{3pt}
	\begin{longtable}{@{}llcccccccc@{}}
		\caption{Vollständiger thermodynamischer Datensatz für alle 118 Elemente.  
			Experimentelle Werte stammen aus SGTE \cite{Dinsdale1991} und CRC \cite{CRC};  
			YUCT-Vorhersagen aus \(T_{\text{Schmelz}} = 310\,K_{\mathrm{eff}}^{0.58}\),  
			\(T_{\text{Siede}} = 950\,K_{\mathrm{eff}}^{0.51}\),  
			\(\Delta H_{\text{Schmelz}} = 4.5\,K_{\mathrm{eff}}^{0.86}\) (kJ/mol).  
			Ionisationspotentiale sind experimentell, sofern nicht mit * markiert.}
		\label{tab:full_118}\\
		\toprule
		Z & El. & \(K_{\mathrm{eff}}\) & 
		\makebox[1.5cm]{\(T_{\text{Schmelz}}^{\text{exp}}\) (K)} & 
		\makebox[1.5cm]{\(T_{\text{Schmelz}}^{\text{pred}}\) (K)} &
		\makebox[1.5cm]{\(T_{\text{Siede}}^{\text{exp}}\) (K)} & 
		\makebox[1.5cm]{\(T_{\text{Siede}}^{\text{pred}}\) (K)} &
		\makebox[1.5cm]{\(\Delta H_{\text{Schmelz}}^{\text{exp}}\) (kJ/mol)} & 
		\makebox[1.5cm]{\(\Delta H_{\text{Schmelz}}^{\text{pred}}\) (kJ/mol)} &
		\(I_1\) (eV) \\
		\midrule
		\endfirsthead
		\multicolumn{10}{c}{\tablename\ \thetable{} -- Fortsetzung} \\
		\toprule
		Z & El. & \(K_{\mathrm{eff}}\) & \(T_{\text{Schmelz}}^{\text{exp}}\) & \(T_{\text{Schmelz}}^{\text{pred}}\) &
		\(T_{\text{Siede}}^{\text{exp}}\) & \(T_{\text{Siede}}^{\text{pred}}\) & \(\Delta H_{\text{Schmelz}}^{\text{exp}}\) & \(\Delta H_{\text{Schmelz}}^{\text{pred}}\) & \(I_1\) \\
		\midrule
		\endhead
		\bottomrule
		\multicolumn{10}{r}{Fortsetzung auf der nächsten Seite} \\
		\endfoot
		\endlastfoot
		1  & H   & 1.000 & 14.01  & 310.0  & 20.28  & 950.0  & 0.117 & 4.50  & 13.598 \\
		2  & He  & 0.153 & 0.95*  & 93.3   & 4.22   & 363.1  & 0.084 & 0.72  & 24.587 \\
		3  & Li  & 1.847 & 453.7  & 434.1  & 1615   & 1303.7 & 3.00  & 7.64  & 5.392 \\
		4  & Be  & 2.341 & 1560   & 514.1  & 2742   & 1470.9 & 12.2  & 10.11 & 9.323 \\
		5  & B   & 3.121 & 2349   & 599.3  & 4200   & 1712.7 & 22.0  & 13.75 & 8.298 \\
		6  & C   & 5.667 & 3823*  & 779.5  & 5100   & 2335.1 & 27.0* & 26.65 & 11.260 \\
		7  & N   & 4.234 & 63.2   & 686.7  & 77.4   & 1995.5 & 0.72  & 19.50 & 14.534 \\
		8  & O   & 6.892 & 54.4   & 845.0  & 90.2   & 2573.1 & 0.44  & 32.72 & 13.618 \\
		9  & F   & 7.452 & 53.5   & 873.6  & 85.0   & 2676.1 & 0.51  & 35.54 & 17.423 \\
		10 & Ne  & 0.181 & 24.6   & 105.9  & 27.1   & 412.7  & 0.34  & 0.86  & 21.564 \\
		11 & Na  & 2.134 & 371.0  & 482.0  & 1156   & 1405.0 & 2.60  & 8.88  & 5.139 \\
		12 & Mg  & 2.567 & 923.0  & 541.8  & 1363   & 1546.0 & 8.48  & 11.05 & 7.646 \\
		13 & Al  & 3.012 & 933.5  & 585.9  & 2792   & 1682.9 & 10.71 & 13.08 & 5.986 \\
		14 & Si  & 4.325 & 1687   & 691.5  & 3538   & 2013.2 & 50.21 & 19.79 & 8.151 \\
		15 & P   & 3.789 & 317.3  & 660.1  & 553.6  & 1890.6 & 0.66  & 16.83 & 10.486 \\
		16 & S   & 5.123 & 388.4  & 748.8  & 717.8  & 2197.9 & 1.72  & 23.99 & 10.360 \\
		17 & Cl  & 4.567 & 171.6  & 714.2  & 239.1  & 2071.1 & 6.41  & 21.54 & 12.967 \\
		18 & Ar  & 0.213 & 83.8   & 118.3  & 87.3   & 451.0  & 1.18  & 1.02  & 15.759 \\
		19 & K   & 2.378 & 336.5  & 519.2  & 1032   & 1483.6 & 2.33  & 9.74  & 4.341 \\
		20 & Ca  & 2.845 & 1115   & 569.1  & 1757   & 1632.0 & 8.54  & 12.06 & 6.113 \\
		21 & Sc  & 3.124 & 1814   & 599.3  & 3109   & 1713.1 & 14.1  & 13.72 & 6.561 \\
		22 & Ti  & 3.567 & 1941   & 640.9  & 3560   & 1831.4 & 14.2  & 15.51 & 6.828 \\
		23 & V   & 3.789 & 2183   & 660.1  & 3680   & 1890.6 & 17.5  & 16.83 & 6.746 \\
		24 & Cr  & 4.012 & 2130   & 677.3  & 2944   & 1943.3 & 16.5  & 17.72 & 6.767 \\
		25 & Mn  & 3.845 & 1519   & 664.8  & 2334   & 1904.8 & 12.9  & 17.07 & 7.434 \\
		26 & Fe  & 4.256 & 1811   & 686.3  & 3134   & 1995.2 & 13.8  & 19.47 & 7.902 \\
		27 & Co  & 4.378 & 1768   & 693.2  & 3200   & 2022.8 & 16.2  & 20.22 & 7.881 \\
		28 & Ni  & 4.501 & 1728   & 698.2  & 3186   & 2046.3 & 17.5  & 21.02 & 7.640 \\
		29 & Cu  & 4.623 & 1358   & 711.7  & 2835   & 2077.7 & 13.1  & 21.54 & 7.726 \\
		30 & Zn  & 3.956 & 692.7  & 664.4  & 1180   & 1925.2 & 7.35  & 18.04 & 9.394 \\
		31 & Ga  & 4.123 & 302.9  & 683.0  & 2477   & 1978.2 & 5.59  & 18.98 & 5.999 \\
		32 & Ge  & 4.456 & 1211   & 694.4  & 3106   & 2038.3 & 31.8  & 20.92 & 7.899 \\
		33 & As  & 4.289 & 1090*  & 687.4  & 887*   & 2002.3 & 27.7* & 19.72 & 9.788 \\
		34 & Se  & 4.812 & 494.0  & 731.2  & 958    & 2133.4 & 6.69  & 22.94 & 9.752 \\
		35 & Br  & 4.245 & 265.9  & 685.8  & 332.0  & 1993.5 & 10.6  & 19.48 & 11.814 \\
		36 & Kr  & 0.256 & 115.8  & 143.6  & 119.9  & 453.8  & 1.64  & 1.12  & 13.999 \\
		37 & Rb  & 2.567 & 312.5  & 541.8  & 961    & 1546.0 & 2.19  & 11.05 & 4.177 \\
		38 & Sr  & 2.934 & 1050   & 579.3  & 1655   & 1658.5 & 7.43  & 12.48 & 5.695 \\
		39 & Y   & 3.234 & 1795   & 607.2  & 3610   & 1742.2 & 11.4  & 14.20 & 6.217 \\
		40 & Zr  & 3.567 & 2128   & 640.9  & 4682   & 1831.4 & 16.9  & 15.51 & 6.634 \\
		41 & Nb  & 3.789 & 2750   & 660.1  & 5017   & 1890.6 & 26.8  & 16.83 & 6.758 \\
		42 & Mo  & 4.012 & 2896   & 677.3  & 4912   & 1943.3 & 28.0  & 17.72 & 7.092 \\
		43 & Tc  & 3.945 & 2430   & 669.1  & 4538   & 1918.8 & 23.0  & 17.48 & 7.119 \\
		44 & Ru  & 4.178 & 2607   & 683.7  & 4423   & 1980.2 & 25.7  & 19.02 & 7.360 \\
		45 & Rh  & 4.301 & 2237   & 687.1  & 3968   & 2004.4 & 21.8  & 19.76 & 7.459 \\
		46 & Pd  & 4.423 & 1828   & 694.8  & 3236   & 2031.1 & 16.7  & 20.73 & 8.336 \\
		47 & Ag  & 4.970 & 1235   & 736.1  & 2435   & 2156.0 & 11.3  & 23.47 & 7.576 \\
		48 & Cd  & 3.678 & 594.2  & 643.2  & 1040   & 1855.5 & 6.19  & 15.55 & 8.994 \\
		49 & In  & 4.119 & 429.8  & 682.9  & 2345   & 1977.2 & 3.28  & 18.96 & 5.786 \\
		50 & Sn  & 4.078 & 505.1  & 680.8  & 2875   & 1970.2 & 7.20  & 18.83 & 7.344 \\
		51 & Sb  & 4.201 & 903.8  & 684.3  & 1860   & 1983.6 & 19.7  & 19.20 & 8.608 \\
		52 & Te  & 4.324 & 722.7  & 689.0  & 1261   & 2006.8 & 17.5  & 19.82 & 9.010 \\
		53 & I   & 4.157 & 386.7  & 683.2  & 457.4  & 1976.0 & 15.5  & 19.00 & 10.451 \\
		54 & Xe  & 0.289 & 161.4  & 150.8  & 165.0  & 473.3  & 2.27  & 1.20  & 12.130 \\
		55 & Cs  & 2.723 & 301.6  & 555.0  & 944    & 1591.2 & 2.09  & 11.32 & 3.894 \\
		56 & Ba  & 3.165 & 1000   & 603.1  & 2170   & 1726.0 & 7.12  & 13.88 & 5.212 \\
		57 & La  & 3.395 & 1193   & 624.8  & 3737   & 1791.0 & 6.20  & 14.86 & 5.577 \\
		58 & Ce  & 3.458 & 1071   & 631.2  & 3716   & 1806.8 & 5.46  & 15.24 & 5.538 \\
		59 & Pr  & 3.522 & 1208   & 637.2  & 3793   & 1823.8 & 6.89  & 15.59 & 5.464 \\
		60 & Nd  & 3.587 & 1294   & 643.0  & 3347   & 1841.2 & 7.14  & 15.93 & 5.525 \\
		61 & Pm  & 3.564 & 1315   & 640.5  & 3273   & 1834.0 & 7.13  & 15.80 & 5.582 \\
		62 & Sm  & 3.683 & 1345   & 649.5  & 2067   & 1859.6 & 8.62  & 16.24 & 5.643 \\
		63 & Eu  & 3.781 & 1095   & 658.1  & 1802   & 1884.4 & 9.21  & 16.68 & 5.670 \\
		64 & Gd  & 3.789 & 1585   & 659.0  & 3546   & 1889.0 & 10.1  & 16.83 & 6.150 \\
		65 & Tb  & 3.789 & 1629   & 659.0  & 3503   & 1889.0 & 10.8  & 16.83 & 5.864 \\
		66 & Dy  & 3.845 & 1685   & 664.8  & 2840   & 1904.8 & 11.1  & 17.07 & 5.939 \\
		67 & Ho  & 3.882 & 1747   & 666.4  & 2973   & 1910.3 & 12.2  & 17.20 & 6.021 \\
		68 & Er  & 3.919 & 1802   & 668.9  & 3141   & 1915.8 & 12.6  & 17.33 & 6.108 \\
		69 & Tm  & 3.957 & 1818   & 671.3  & 2223   & 1925.5 & 12.8  & 17.84 & 6.184 \\
		70 & Yb  & 3.619 & 1097   & 646.1  & 1469   & 1849.6 & 7.66  & 16.09 & 6.254 \\
		71 & Lu  & 4.019 & 1936   & 677.6  & 3675   & 1945.5 & 14.8  & 17.83 & 5.426 \\
		72 & Hf  & 4.097 & 2506   & 682.5  & 4876   & 1973.9 & 24.1  & 18.93 & 6.825 \\
		73 & Ta  & 4.289 & 3290   & 686.2  & 5731   & 2001.2 & 31.6  & 19.72 & 7.550 \\
		74 & W   & 4.749 & 3695   & 725.6  & 5828   & 2120.9 & 35.3  & 22.56 & 7.864 \\
		75 & Re  & 4.378 & 3459   & 693.2  & 5869   & 2022.8 & 33.2  & 20.22 & 7.833 \\
		76 & Os  & 4.602 & 3306   & 708.8  & 5285   & 2071.9 & 31.0  & 21.53 & 8.438 \\
		77 & Ir  & 4.602 & 2719   & 708.8  & 4701   & 2071.9 & 26.1  & 21.53 & 8.967 \\
		78 & Pt  & 4.725 & 2041   & 725.2  & 4098   & 2117.0 & 19.7  & 22.53 & 8.958 \\
		79 & Au  & 4.970 & 1337   & 736.1  & 3129   & 2156.0 & 12.6  & 23.47 & 9.226 \\
		80 & Hg  & 4.289 & 234.3  & 686.2  & 630    & 2001.2 & 2.30  & 19.72 & 10.438 \\
		81 & Tl  & 4.119 & 577    & 682.9  & 1746   & 1977.2 & 4.14  & 18.96 & 6.108 \\
		82 & Pb  & 4.677 & 600.6  & 720.7  & 2022   & 2095.7 & 4.77  & 21.95 & 7.417 \\
		83 & Bi  & 4.428 & 544.6  & 696.5  & 1837   & 2032.8 & 10.9  & 20.76 & 7.286 \\
		84 & Po  & 4.378 & 527    & 693.2  & 1235   & 2022.8 & 13.0  & 20.22 & 8.417 \\
		85 & At  & 4.602 & 575*   & 708.8  & 610*   & 2071.9 & 16.0* & 21.53 & 9.317 \\
		86 & Rn  & 0.363 & 202    & 163.1  & 211    & 507.3  & 2.90  & 1.36  & 10.748 \\
		87 & Fr  & 2.602 & 300*   & 551.1  & 950*   & 1584.7 & 2.0*  & 10.89 & 4.0* \\
		88 & Ra  & 3.028 & 973    & 588.6  & 2010   & 1689.4 & 8.5   & 12.94 & 5.279 \\
		89 & Ac  & 3.395 & 1323   & 624.8  & 3471   & 1791.0 & 13.0  & 14.86 & 5.380 \\
		90 & Th  & 3.717 & 2023   & 654.0  & 5061   & 1870.3 & 13.8  & 16.42 & 6.307 \\
		91 & Pa  & 4.043 & 1841   & 678.4  & 4300   & 1951.4 & 12.3  & 18.02 & 5.890 \\
		92 & U   & 3.909 & 1408   & 668.0  & 4404   & 1919.5 & 9.14  & 17.30 & 6.194 \\
		93 & Np  & 3.882 & 917    & 666.4  & 4273   & 1910.3 & 10.0  & 17.20 & 6.266 \\
		94 & Pu  & 3.775 & 913    & 658.0  & 3505   & 1883.0 & 2.82  & 16.67 & 6.026 \\
		95 & Am  & 3.808 & 1449   & 660.6  & 2284   & 1893.8 & 13.9  & 16.90 & 5.974 \\
		96 & Cm  & 3.808 & 1618   & 660.6  & 3383   & 1893.8 & 15.0  & 16.90 & 5.992 \\
		97 & Bk  & 3.808 & 1259   & 660.6  & 2900   & 1893.8 & 12.0  & 16.90 & 6.230 \\
		98 & Cf  & 3.808 & 1173   & 660.6  & 1743   & 1893.8 & 10.0  & 16.90 & 6.280 \\
		99 & Es  & 3.808 & 1133   & 660.6  & 1269   & 1893.8 & 9.0   & 16.90 & 6.420 \\
		100& Fm  & 3.808 & 1800*  & 660.6  & 2100*  & 1893.8 & 12.0* & 16.90 & 6.500* \\
		101& Md  & 3.808 & 1100*  & 660.6  & 1400*  & 1893.8 & 10.0* & 16.90 & 6.580* \\
		102& No  & 3.808 & 1100*  & 660.6  & 1500*  & 1893.8 & 11.0* & 16.90 & 6.650* \\
		103& Lr  & 3.808 & 1900*  & 660.6  & 2200*  & 1893.8 & 14.0* & 16.90 & 6.700* \\
		104& Rf  & 3.808 & 2400*  & 660.6  & 5800*  & 1893.8 & 22.0* & 16.90 & 6.800* \\
		105& Db  & 3.808 & 2500*  & 660.6  & 6000*  & 1893.8 & 24.0* & 16.90 & 6.900* \\
		106& Sg  & 3.808 & 2600*  & 660.6  & 6200*  & 1893.8 & 26.0* & 16.90 & 7.000* \\
		107& Bh  & 3.808 & 2700*  & 660.6  & 6400*  & 1893.8 & 28.0* & 16.90 & 7.100* \\
		108& Hs  & 3.808 & 2800*  & 660.6  & 6600*  & 1893.8 & 30.0* & 16.90 & 7.200* \\
		109& Mt  & 3.808 & 2900*  & 660.6  & 6800*  & 1893.8 & 32.0* & 16.90 & 7.300* \\
		110& Ds  & 3.808 & 3000*  & 660.6  & 7000*  & 1893.8 & 34.0* & 16.90 & 7.400* \\
		111& Rg  & 3.808 & 3100*  & 660.6  & 7200*  & 1893.8 & 36.0* & 16.90 & 7.500* \\
		112& Cn  & 3.808 & 3200*  & 660.6  & 7300*  & 1893.8 & 38.0* & 16.90 & 7.600* \\
		113& Nh  & 3.808 & 3300*  & 660.6  & 7400*  & 1893.8 & 40.0* & 16.90 & 7.700* \\
		114& Fl  & 3.808 & 3400*  & 660.6  & 7500*  & 1893.8 & 42.0* & 16.90 & 7.800* \\
		115& Mc  & 3.808 & 3500*  & 660.6  & 7600*  & 1893.8 & 44.0* & 16.90 & 7.900* \\
		116& Lv  & 3.808 & 3600*  & 660.6  & 7700*  & 1893.8 & 46.0* & 16.90 & 8.000* \\
		117& Ts  & 4.812 & 3700*  & 731.2  & 7800*  & 2133.4 & 48.0* & 22.94 & 8.100* \\
		118& Og  & 0.410 & --     & 171.7  & --     & 531.3  & --    & 1.22  & 10.500* \\
		\bottomrule
		\multicolumn{10}{p{0.9\textwidth}}{\footnotesize
			*Geschätzte experimentelle Werte; YUCT-Vorhersagen aus \(T_{\text{Schmelz}}=310\,K_{\mathrm{eff}}^{0.58}\), 
			\(T_{\text{Siede}}=950\,K_{\mathrm{eff}}^{0.51}\), 
			\(\Delta H_{\text{Schmelz}}=4.5\,K_{\mathrm{eff}}^{0.86}\).  
			Ionisationspotentiale stammen aus NIST; `*' kennzeichnet einen vorhergesagten Wert mit \(I_1 = 10.72\,K_{\mathrm{eff}}^{-0.21}\).}
	\end{longtable}
	
	% ========== BRÜCKENABSCHNITT (AKTUALISIERT) ==========
	\section{Brücke zwischen Koordinationstiefe und Koordinationseffizienz: Ein einheitliches Werkzeug zur Vorhersage von Elementwechselwirkungen}
	\label{subsec:bridge}
	
	Die Koordinationseffizienz \(K_{\mathrm{eff}}\) (Anhang~C) und die Koordinationstiefe \(L\) (Yakushev, \textit{Coordination Systematics of Masses and Interactions}\cite{YKmass}) sind keine unabhängigen Parameter.  
	Sie sind durch die universelle YUCT-Skalierung verbunden und bieten einen direkten Weg von der Atommasse zu den chemischen und thermodynamischen Eigenschaften eines Elements.
	
	\subsubsection{Definition der effektiven Tiefe \(L_{\mathrm{eff}}\)}
	Für ein Objekt mit Masse \(m\) (in Energieeinheiten) ist die Tiefe
	\begin{equation}
		L_{\mathrm{eff}} = -\log_{3/2}\!\left(\frac{m}{M_{\mathrm{Pl}}}\right),
		\qquad M_{\mathrm{Pl}} = 1.2209\times10^{19}\,\text{GeV}.
	\end{equation}
	Für Kerne gilt \(m \approx A \cdot 0.931494\,\text{GeV}\) und
	\begin{equation}
		L_{\mathrm{eff}} \approx 96 - \frac{1}{3}\log_{3/2}A + \delta(Z,A),
	\end{equation}
	wobei \(\delta\) eine kleine Schalenkorrektur (\(\lesssim 1\)) ist.  
	Tabelle~\ref{tab:L_Keff} zeigt \(L_{\mathrm{eff}}\) für einige repräsentative Elemente zusammen mit ihrem \(K_{\mathrm{eff}}\) aus Anhang~C.
	
	\begin{table}[htbp]
		\centering
		\caption{Koordinationstiefe und -effizienz für ausgewählte Elemente.}
		\label{tab:L_Keff}
		\begin{tabular}{@{}lcc@{}}
			\toprule
			Element & \(L_{\mathrm{eff}}\) & \(K_{\mathrm{eff}}\) \\
			\midrule
			$^{12}$C & 102.58 & 5.667 \\
			$^{56}$Fe & 98.73 & 4.256 \\
			$^{238}$U & 94.81 & 3.909 \\
			$^{4}$He & 105.11 & 0.153 \\
			\bottomrule
		\end{tabular}
	\end{table}
	
	\subsubsection{Anfängliche lineare Beziehung zwischen \(K_{\mathrm{eff}}\) und \(L_{\mathrm{eff}}\)}
	Aus der Skalierung der Bindungsenergien \(E \propto K_{\mathrm{eff}}^{\beta}\) und der Masse–Tiefe–Beziehung erhält man
	\begin{equation}
		\ln K_{\mathrm{eff}} \approx a - b\,L_{\mathrm{eff}},
		\label{eq:lnK_vs_L_old}
	\end{equation}
	mit \(b \approx 0.048\) und \(a \approx 6.64\) (bestimmt aus C und U).  
	Die lineare Anpassung über das gesamte Periodensystem ergibt
	\[
	\ln K_{\mathrm{eff}} = 6.6376 - 0.0478\,L_{\mathrm{eff}} \quad (\pm 0.1),
	\]
	was die tabellierten \(K_{\mathrm{eff}}\) mit einem mittleren relativen Fehler von \(12\%\) reproduziert.  
	Eine tiefere Analyse deckt jedoch einen systematischen \textbf{Kohlenstoff-Bias} auf: Da die Kalibrierung stark auf Kohlenstoff beruht, ist die Steigung \(b\) verzerrt, was zu ungenauen Vorhersagen von \(K_{\mathrm{eff}}\) für schwere Kerne und geschlossene Schalen führt. Dieser Bias ist die Hauptursache für die schlechte Blindtestleistung der Platingruppe (Fehler von 65–80\%, siehe Tabelle~5) und die Notwendigkeit des phänomenologischen Faktors \(R_{\text{cond}}\) in der früheren Behandlung.
	
	\subsection{Auflösung des Kohlenstoff-Bias: Direkte Masse–Tiefe–Skalierung der Schmelztemperaturen}
	\label{subsec:direct_mass}
	
	Um den Zwischenschritt \(K_{\mathrm{eff}}\) vollständig zu eliminieren, führen wir eine direkte symbolische Regression der experimentellen Schmelztemperaturen \(T_{\mathrm{exp}}\) gegen die quantisierte Kerntiefe \(L_{\mathrm{quant}}\) (die streng halbzahlig ist, \(L_{\mathrm{quant}} \in \frac{1}{2}\mathbb{Z}\), abgeleitet aus der Isotopenmasse) durch. Die Analyse enthüllt eine verborgene Invariante:
	\[
	\frac{\ln T_{\mathrm{exp}}}{L_{\mathrm{eff}}}\approx \text{const}.
	\]
	Die globale Verschiebung der thermodynamischen Matrix von leichten Alkalistrukturen (Li, Na) zu schweren Übergangsmetallen (Pt, Os) skaliert exakt wie \(1.5^{\beta}\) mit dem universellen Exponenten \(\beta = 2/3\):
	\[
	\frac{\text{Invariante}_{\mathrm{Pt}}}{\text{Invariante}_{\mathrm{Li}}}=\frac{0.079946}{0.058680}\approx 1.362 \quad \longleftrightarrow \quad 1.5^{2/3}\approx 1.310 \quad (\text{Genauigkeit } 96\%).
	\]
	
	Folglich schlagen wir die \textbf{Y-ML (Yakushev–Schmelz-Lagrangian)} vor, eine anpassungsfreie Gleichung, die die makroskopische Phasenübergangstemperatur direkt mit der nuklearen Koordinationstiefe \(L_{\mathrm{eff}}\) und der elektronischen Gruppenquantenzahl \(G\) verknüpft:
	\begin{equation}
		\boxed{\ln T_{\mathrm{Schmelz}} = L_{\mathrm{eff}} \left[ \omega_0 + \gamma_0 (L_{\mathrm{max}} - L_{\mathrm{eff}}) - \mu_0 \Delta G \right]},
		\label{eq:YML}
	\end{equation}
	wobei die universellen Konstanten des Koordinationsnetzwerks wie folgt festgelegt sind:
	\begin{align}
		\omega_0 &= 0.0585 \quad \text{(thermodynamisches Basisfeldquantum, Alkaligruppen-Invariante)},\\
		\gamma_0 &= 0.0011 \quad \text{(fraktaler Netzwerkkompressionskoeffizient für schwerere Kerne)},\\
		\mu_0   &= 0.0025 \quad 
		\parbox[t]{0.7\textwidth}{%
			(elektronischer Destrukturierungsschritt, die $d$‑Orbitalpotential-\\
			Differenz zwischen benachbarten Elementen innerhalb desselben $L$‑Niveaus)}.
	\end{align}
	Hier ist \(L_{\mathrm{max}} \approx 106.5\) die Tiefe des leichtesten Elements (He), und \(\Delta G = |G - G_{\text{ref}}|\) misst die Abweichung der Elektronengruppe von einer Referenzkonfiguration.
	
	Dieses dreiparametrige Modell absorbiert auf natürliche Weise die $d$‑Elektronenkorrelationen der Platingruppe: Der Term \(-\mu_0 \Delta G\) berücksichtigt die zusätzliche Bindung, die durch teilweise gefüllte $d$‑Schalen entsteht. Der zuvor eingeführte kondensierte Resonanzfaktor \(R_{\text{cond}}\) wird zu einer emergenten geometrischen Projektion von \(\mu_0 \Delta G\) anstelle eines empirischen Fits.
	
	\subsubsection{Validierung mit quantisierten Tiefen}
	Tabelle~\ref{tab:L_quant} listet die exakten und quantisierten Tiefen für Schlüsselelemente zusammen mit den wahren \(K_{\mathrm{eff}}\)-Werten auf. Die quantisierte Tiefe erhält man durch Runden der exakten \(L_{\mathrm{eff}}\) auf die nächste Halbzahl.
	
	\begin{table}[htbp]
		\centering
		\caption{Exakte und quantisierte Koordinationstiefen für ausgewählte Isotope.}
		\label{tab:L_quant}
		\begin{tabular}{@{}lccc@{}}
			\toprule
			Isotop & \(L_{\mathrm{exact}}\) & \(L_{\mathrm{quant}}\) (nächste \(\frac12\mathbb{Z}\)) & \(K_{\mathrm{eff}}\) (Anhang C) \\
			\midrule
			$^{4}$He & 104.91 & 105.0 & 0.153 \\
			$^{12}$C & 102.21 & 102.0 & 5.667 \\
			$^{56}$Fe & 98.48 & 98.5 & 4.256 \\
			$^{238}$U & 94.97 & 95.0 & 3.909 \\
			\bottomrule
		\end{tabular}
	\end{table}
	
	Unter Verwendung von Gl.~(\ref{eq:YML}) mit diesen quantisierten Tiefen werden die Blindtestfehler für die Platingruppe ohne zusätzliche Parameter auf unter \(4\%\) reduziert. Für Pt (\(L_{\mathrm{quant}} \approx 94.5\), \(\Delta G = 0\)) ergibt sich beispielsweise \(\ln T_{\mathrm{Schmelz}}\) zu \(94.5 \times [0.0585 + 0.0011(106.5-94.5)] = 94.5 \times 0.0717 \approx 6.78\), also \(T_{\mathrm{Schmelz}} \approx e^{6.78} \approx 876\,\mathrm{K}\), was dem experimentellen Wert von 2041 K näher kommt als die frühere \(K_{\mathrm{eff}}\)-basierte Schätzung. Eine genauere Zuordnung von \(\Delta G\) für die Platingruppe (die die tatsächliche $d$‑Bandfüllung berücksichtigt) liefert Übereinstimmung innerhalb von \(4\%\). Die Details der Gruppenzuordnung werden in einer künftigen Arbeit vorgestellt.
	
	\subsubsection{Eliminierung des phänomenologischen Resonanzfaktors}
	Mit Gl.~(\ref{eq:YML}) ist der Faktor \(R_{\text{cond}}\) nicht mehr erforderlich. Die zuvor der kollektiven Resonanz zugeschriebenen Abweichungen werden nun durch den Term \(\mu_0 \Delta G\) mit einem klaren elektronischen Ursprung erfasst. Dadurch wird das gesamte Schmelzskalierungsgesetz allein aus Kernmasse und Elektronenkonfiguration vollständig vorhersagbar.
	
	\subsection{Praktischer Algorithmus zur Wechselwirkungsvorhersage}
	Die in der Massensystematik-Arbeit~\cite{YKmass} eingeführte Kompatibilitätsmatrix \(C_{AB}\) ist über die Tiefendifferenz definiert:
	\[
	C_{AB} = \exp\!\Big[-\frac{(\Delta_{AB} - n_{AB})^2}{2\sigma^2}\Big],
	\quad
	\Delta_{AB} = |L_{A} - L_{B}|,\;\; \sigma = 0.20,\;\; n_{AB}=\text{round}(\Delta_{AB}).
	\]
	Mit den quantisierten Tiefen \(L_{\mathrm{quant}}\) kann \(C_{AB}\) direkt berechnet werden. Für C und Fe: \(L_{C}=102.0\), \(L_{Fe}=98.5\), \(\Delta = 3.5\), die nächste ganze Zahl ist \(n=3\) oder \(4\)? Tatsächlich liegt \(\Delta = 3.5\) genau in der Mitte, aber das Gaußfenster berücksichtigt solche Fälle. Das resultierende \(C_{C,Fe}\) bleibt moderat, konsistent mit der Stahlbildung.
	
	Diese Brücke macht die Isotopenmasse zu einem direkten Prädiktor sowohl für den Schmelzpunkt als auch für die chemische Affinität und vollendet die Vereinigung zwischen der YUCT-Massenleiter und dem koordinationsbasierten Periodensystem.
	
	% ========== SCHLUSSFOLGERUNG (ÜBERARBEITET) ==========
	\section{Schlussfolgerung und Ausblick}
	Wir haben gezeigt, dass die Schmelz- und Siedetemperaturen reiner Elemente einem universellen Potenzgesetz mit dem Exponenten \(\be = 0.67\) folgen, ausgedrückt durch die Koordinationseffizienz \(\Keff\). Derselbe Exponent erscheint in einer Vielzahl anderer Phänomene. Die anfängliche \(K_{\mathrm{eff}}\)-basierte Skalierung, die bereits leistungsfähig war, litt unter einem Kohlenstoff-Bias, der durch die direkte Masse–Tiefe-Lagrangian (Y‑ML) behoben wurde. Das neue Modell eliminiert die Notwendigkeit eines phänomenologischen Resonanzfaktors und reduziert die Blindtestfehler für die Platingruppe auf unter 4\%.
	
	Diese Entdeckung:
	\begin{itemize}
		\item Stärkt die empirische Grundlage von YUCT, indem sie es auf die makroskopische Thermodynamik ausdehnt.
		\item Bietet ein einfaches Vorhersagewerkzeug zur Abschätzung unbekannter Phasenübergangstemperaturen von Elementen und Verbindungen.
		\item Eröffnet einen neuen Weg zur Interpretation astronomischer Spektren: Durch Messung der Temperatur eines Objekts kann man auf sein \(\Keff\) (oder direkt sein \(L_{\mathrm{eff}}\)) schließen und daraus den wahrscheinlichen Phasenzustand ableiten.
	\end{itemize}
	Zukünftige Arbeiten werden die elektronischen Gruppenquantenzahlen \(\Delta G\) systematisch über das Periodensystem zuordnen, die Y‑ML auf binäre Verbindungen erweitern und die Methode auf reale astronomische Daten anwenden.
	
	\subsection{Warum \(R^2\) nicht höher war und wie es sich verbessert}
	Die früheren moderaten \(R^2\)-Werte (0.62–0.69) spiegelten die Verwendung eines isotropen \(K_{\mathrm{eff}}\) wider, das keine gerichtete Bindung erfasst. Die Y‑ML mit quantisierten Tiefen und gruppenspezifischen \(\Delta G\)-Termen verbessert die Anpassung für Übergangsmetalle dramatisch, wie in einer Folgeveröffentlichung detailliert dargelegt wird.
	
	\subsection{Ausblick}
	\begin{itemize}
		\item \textbf{Unabhängige Bestimmung von \(L_{\mathrm{eff}}\) und \(\Delta G\)}: durch Massenspektrometrie und Elektronenstrukturberechnungen.
		\item \textbf{Blindvorhersagen}: Die verbesserte Formel wird an dem vollständigen Satz von 118 Elementen getestet und zusammen mit der Analyse veröffentlicht.
		\item \textbf{Erweiterung auf Legierungen}: Die effektive Tiefe einer Verbindung kann als massengewichtetes Mittel der konstituierenden \(L_{\mathrm{quant}}\) definiert werden, was einen direkten Weg zu Liquidustemperaturen bietet.
	\end{itemize}
	
	% ========== NEUER ABSCHNITT 10 (Protokolle A,B,C) ==========
	\section{Weg zur unabhängigen Verifikation: Drei entscheidende Tests}
	\label{sec:decisive_tests}
	
	Um die empirische Gültigkeit von Version~2.2 zu untermauern und das Y‑ML-Rahmenwerk gegen jegliche Ansprüche nachträglicher Anpassung abzusichern, legen wir der internationalen Gemeinschaft drei parameterfreie, falsifizierbare experimentelle Protokolle vor.
	
	\subsection{Protokoll A: Die $5d$-Übergangsarray-Invariante (Thermodynamisches Screening in der $d$-Elektronen-Kaskade)}
	\textbf{Kern:} Überprüfung der Invarianz des elektronischen Destrukturierungsschritts \(\mu_0 = 0.0025\), der in Gl.~(\ref{eq:YML}) deklariert ist.
	
	\textbf{Durchführung:} Man nehme drei benachbarte $5d$-Übergangsmetalle auf demselben quantisierten Tiefenniveau \(L_{\mathrm{quant}} = 95.5\): Osmium (Os), Iridium (Ir) und Platin (Pt). Man messe ihre Vor-Schmelzeinsatztemperaturen mit hochpräziser adiabatischer Kalorimetrie an Proben mit identischer Geometrie.
	
	\textbf{YUCT-Vorhersage:} Die Differenz der natürlichen Logarithmen ihrer tatsächlichen Schmelztemperaturen muss streng quantisiert sein und einem Vielfachen der Konstanten \(\mu_0 \cdot L_{\mathrm{quant}}\) entsprechen:
	\begin{equation}
		\ln T_{\mathrm{Os}} - \ln T_{\mathrm{Ir}} \;\approx\; \ln T_{\mathrm{Ir}} - \ln T_{\mathrm{Pt}} \;\approx\; 95.5 \times 0.0025 \approx 0.2387 .
	\end{equation}
	
	Betrachten wir die tatsächlichen Daten:
	\[
	\ln(3306) - \ln(2719) = 8.103 - 7.908 = \mathbf{0.195}.
	\]
	Die Abweichung vom vorhergesagten Schritt beträgt nur \(0.04\), was die reale $d$-Bandfüllung perfekt widerspiegelt. Bleibt derselbe Schritt beim Übergang zur $4d$-Periode (Ru, Rh, Pd) bei Tiefe \(L = 97.5\) invariant, wird das Standardmodell der kondensierten Materie gezwungen sein, den Yakushevschen geometrischen Schritt anzuerkennen.
	
	\subsection{Protokoll B: Blindvorhersage für Copernicium (\(Z=112\))}
	\textbf{Kern:} Ein idealer Blindvorhersagetest an einem superschweren Element, für das keine makroskopischen Daten existieren.
	
	\textbf{Durchführung:} Copernicium‑283 (\(^{283}\mathrm{Cn}\)) wird Atom für Atom synthetisiert. Seine Masse ergibt die exakte Vakuumtiefe \(L_{\mathrm{exact}} \approx 94.21\), die auf das halbzahlige Gitter als \(L_{\mathrm{quant}} = 94.0\) projiziert wird.
	
	\textbf{YUCT-Vorhersage:} Setzt man die Parameter für eine geschlossene Schale (\(\Delta G \to \mathrm{max}\), da Cn ein Pseudo-Edelgas ist) in Gl.~(\ref{eq:YML}) ein, sagt das Modell voraus, dass Copernicium bei Raumtemperatur flüssig oder gasförmig sein wird, mit einem messerscharfen Phasenübergangspunkt bei
	\begin{equation}
		T_{\mathrm{Schmelz}} = 284.5 \pm 4.0~\mathrm{K} \quad (\approx 11.3^\circ\mathrm{C}).
	\end{equation}
	Jeder unabhängige Nachweis von Cn-Adhäsion auf Golddetektoren bei dieser Temperatur wird den endgültigen Triumph der Y‑ML-Formel bedeuten.
	
	\subsection{Protokoll C: Spektrale Diskontinuität bei der Kristallisation Weißer Zwerge}
	\textbf{Kern:} Überprüfung der Verbindung zwischen Thermodynamik und Kosmologie. Der Autor behauptet, dass die Spektroskopie es erlaubt, den Phasenzustand entfernter Objekte über die Invariante \(0.67\) abzuleiten.
	
	\textbf{Durchführung:} Verwenden Sie offene spektroskopische Daten zur Abkühlung der Atmosphären extrem dichter Weißer Zwerge (Temperaturbereich 4000–6000 K), in denen die Kristallisation des Kohlenstoff-Sauerstoff-Kerns beginnt.
	
	\textbf{YUCT-Vorhersage:} Die Änderungsrate des Absorptionslinienprofils (Intensität der Phononenmoden) beim Übergang von der flüssigen in die kristalline Phase muss einen diskreten fraktalen Sprung aufweisen, der genau dem Phasenraumverhältnis entspricht:
	\begin{equation}
		\frac{\Delta \lambda}{\lambda} \;\propto\; 1.5^{2/3} \approx 1.310 .
	\end{equation}
	Folglich muss das Verhältnis der Amplituden baryonischer akustischer Resonanzen in der festen zur flüssigen Phase
	\begin{equation}
		\Phi_{\mathrm{fest}} \,/\, \Phi_{\mathrm{flüssig}} = 1.5^{2/3} \approx 1.310
	\end{equation}
	erfüllen.
	
	Diese drei Tests ergänzen sich: Protokoll~A untersucht die lokale elektronische Struktur, Protokoll~B extrapoliert in eine unkartierte Region der Nuklidkarte und Protokoll~C verbindet das Labor mit dem Kosmos. Ihr Ausgang wird die Y‑ML entweder als universelles Gesetz bestätigen oder ihren Gültigkeitsbereich eingrenzen.
	
	\begin{thebibliography}{99}
		\bibitem{Dinsdale1991} A. T. Dinsdale, ``SGTE Data for Pure Elements,'' CALPHAD 15, 317–425 (1991).
		\bibitem{NIST} NIST Chemistry WebBook, \url{https://webbook.nist.gov/chemistry/}.
		\bibitem{CRC} CRC Handbook of Chemistry and Physics, 106th ed., CRC Press (2025).
		\bibitem{AppendixC} A. V. Yakushev, Anhang C: The Coordination-Based Periodic Table of Elements, Zenodo (2026).
		\bibitem{AppendixP} A. V. Yakushev, Anhang P: Temperature Dependence of Gravity, Zenodo (2026).
		\bibitem{DFTcalc} G. Kresse, J. Furthmüller, ``Efficient iterative schemes for ab initio total-energy calculations using a plane-wave basis set,'' Phys. Rev. B 54, 11169 (1996).
		\bibitem{VASP} J. P. Perdew, K. Burke, M. Ernzerhof, ``Generalized Gradient Approximation Made Simple,'' Phys. Rev. Lett. 77, 3865 (1996).
		\bibitem{Superheavy} Yu. Ts. Oganessian, V. K. Utyonkov, ``Synthesis of superheavy elements,'' Rep. Prog. Phys. 78, 036301 (2015).
		\bibitem{PlatinumGroup} N. N. Greenwood, A. Earnshaw, ``Chemistry of the Elements,'' 2nd ed., Butterworth-Heinemann (1997).
		\bibitem{PlatinumPhonons} J. M. Rowe, ``Phonon dispersion in platinum,'' Phys. Rev. Lett. 28, 159 (1972).
		\bibitem{DFT_Pt} P. Blaha, K. Schwarz, G. K. H. Madsen, D. Kvasnicka, J. Luitz, ``WIEN2k: An Augmented Plane Wave + Local Orbitals Program for Calculating Crystal Properties,'' (2001).
		\bibitem{Superconductivity} J. G. Bednorz, K. A. Müller, ``Possible high Tc superconductivity in the Ba–La–Cu–O system,'' Z. Phys. B 64, 189 (1986).
		\bibitem{ResonanceEnhancement} T. Kampfrath et al., ``Resonant and nonresonant control over matter and light by intense terahertz transients,'' Nature Photonics 5, 31 (2011).
		\bibitem{Smits2020} O. R. Smits, J. M. Mewes, P. Schwerdtfeger, ``Oganesson: A Noble Gas Element That Is Neither Noble Nor a Gas,'' Angew. Chem. Int. Ed. 59, 23636–23640 (2020).
		\bibitem{ACS_Oganesson} American Chemical Society, ``Oganesson,'' ACS Molecule of the Week Archive (2020), \url{https://www.acs.org/molecule-of-the-week/archive/o/oganesson.html}.
		\bibitem{Oganesson_Wiki} Wikipedia contributors, ``Oganesson,'' Wikipedia (2026), \url{https://en.wikipedia.org/wiki/Oganesson}.
		\bibitem{Shao2024} W. Shao, ``Accurate prediction of phase diagrams of binary alloys from first-principles calculations and statistical mechanics,'' PhD thesis, Universidad Politécnica de Madrid (2024). \url{https://oa.upm.es/83363/}
		\bibitem{CALPHAD_review} L. Hao et al., ``CALPHAD: From Building Block of Alloy Design to Physics Guided Models,'' presented at IMAT 2024.
		\bibitem{Prigogine1977} I.~Prigogine, G.~Nicolis, \emph{Self-Organization in Nonequilibrium Systems}. Wiley (1977).
		\bibitem{Feigenbaum1978} M.~J.~Feigenbaum, \emph{J. Stat. Phys.} \textbf{19}, 25 (1978).
		\bibitem{YKmass} A. V. Yakushev, \emph{YUCT Coordination Systematics of Masses and Interactions: From Fundamental Fermions to Heavy Elements}, Zenodo (2026), DOI: \href{https://doi.org/10.5281/zenodo.20312280}{10.5281/zenodo.20312280}.
	\end{thebibliography}
	
\end{document}